\newpage
\chapter{Conclusion}
\mbox{}\vspace{-\baselineskip}

This report presents a thorough and comprehensive investigation of possibilities to measure color transparency effects with the SPD detector at the NICA collider. As discussed in Ch.~\ref{Chapter_observables}, several observables were found to be sensitive to CT effects, including some polarization observables. This particular study mostly focused on the cross section measurements for the $pd\rightarrow ppn$ process occurring in $d-d$ collisions, and specifically, on the cross section dependence on the transverse momentum of the spectator neutron ($p_{st}$), which was found to be is very informative with respect to CT effects.

In general, the procedure of CT measurements should include two major steps, i.e., (i) experimental measurements of observables sensitive to CT effects and (ii) the subsequent model analysis of the measured observables to pursue estimations of the CT contribution. The latter will be performed in the framework of the generalized eikonal approach supplemented by the quantum diffusion model (QDM) of CT~\cite{Frankfurt:1996uz,PhysRevC.107.014605,Larionov:2025equ,Farrar:1988me}.



According to the available model calculations~\cite{Frankfurt:1996uz,PhysRevC.107.014605,Larionov:2025equ,Farrar:1988me}, the strength of the CT contribution depends on the value of the momentum transferred, i.e., it increases as $Q^2$ grows, and maximal CT contribution is achieved at the $Q^{2}$ values that correspond to $\theta_{\textrm{c.m.}}=90^{o}$ (for a fixed value of $\sqrt{s_{34}}\simeq \frac{\sqrt{s}}{2}$). To facilitate the model analysis, observables sensitive to CT effects should be measured at the highest possible $\theta_{\textrm{c.m.}}$ values. However, taking into account a steep cross section drop with $t$, measurements are proposed to be performed at $\theta_{\textrm{c.m.}}<90^{o}$, and depending on actual experimental conditions, a compromise between the statistics and the $\theta_{\textrm{c.m.}}$ value should be observed.

The performed event rate estimation presented in Ch.~\ref{chapter_rate} proves the overall feasibility of the proposed color transparency measurements at the energy and luminosity values that are expected to be available in the SPD experiment at NICA. According to this estimation, the expected statistics should allow for measurements of the $p_{st}$ dependence of the reaction cross section with reasonable accuracy. 

The measurements are proposed to be performed during the first stage of the NICA collider operation, and, as described in Ch.~\ref{chap:event_selection}, they may encounter a set of experimental challenges. First, in the absence of the TOF system during the first stage, the proton identification will rely mostly on the information from the Straw system, which will be able to adequately perform the $\pi$/K/p separation only for the particle momenta of $\lesssim$~1.5~GeV/c. For higher momentum values, reliable separation of charged hadrons will be challenging, which will put limitations on the achievable values of the coherent length (and therefore, on the strength of the measurable CT contribution).

Second, in the absence of a proper hadron calorimeter in the SPD setup, neutron detection will be performed in the Range system and in the ZDCs.  However, very limited geometric acceptance of the ZDCs will prevent the majority of the forward/backward going neutrons from detection. In addition to that, geometrical acceptance of the Range system will impose limitations on the minimal achievable $p_{st}$ value.

Furthermore, some challenges are also foreseen to emerge on the data analysis stage. One of them will stem from the fact that the Fermi momentum of the proton in a deuteron will not be experimentally accessible, which will result in the Fermi-smearing of the calculated kinematic variables (and the cross sections binned in them). This is a standard situation for deuteron experiments, which will require careful event selection and possible modifications to the theoretical model. Another challenge will be related to the need to select events in which one of the neutrons remained quasi-free, whereas for the other neutron, both quasi-free events and events with soft rescatterings should be included in the analysis. This is, however, a very challenging task as kinematic distributions of both neutrons involved in the process will be indistinguishable from each other. A universal way to approach this challenge in deuteron studies is to utilize missing quantities~\cite{CLAS:2023mfc,my_an_note:2020,my_thesis:2021,Skorodumina:2015rea, CLAS:2022kta,note_mm_distr}, but in this particular experiment, this approach can be complicated by non-beneficial energy-momentum resolution of the experiment. In any case, a detailed and comprehensive examination of different missing quantities (and their components) should be endeavored during the data analysis. 


The performed study allows us to conclude that the proposed CT measurements are anticipated to be rather challenging but still feasible, as many of the projected complications are common for deuteron experiments and, therefore, reliable solutions for them have already become available based on previous studies~\cite{CLAS:2023mfc,my_an_note:2020,my_thesis:2021,Skorodumina:2015rea, CLAS:2022kta,note_mm_distr}.

Also note that, in general, CT measurements benefit from higher collision energies as they deliver larger values of the coherence length and, consequently, a larger overall CT contribution. For this reason, in case of unfavorable conditions during the first stage of the NICA collider operation, the measurements can also be attempted during the second stage, for which larger collision energies and luminosity values will be available along with a complete SPD setup. 

Once a reliable detector simulation and event reconstruction software for the SPD detector is available, this study should be complemented by the examination of reconstructed events and background estimations, which will provide further insights into feasibility of the proposed measurements. Also, for further studies, a more elaborated inquest on polarization observables is required as well as efforts towards the preparation of the detailed run plan.





